\begin{definition}
For a real number $p$ put
\[
  \kappa(p)=\frac16\Biggl(7+\frac1{p-1}+\sqrt{1+\frac{14}{p-1}+\frac1{(p-1)^2}}\,\Biggr).
\]
This is exactly the right-hand side of equation (9) of the paper, that is, the
number which the integer $k\ge1$ and the exponent $\alpha\in(0,1]$ of
Theorem~1 are required to satisfy $k+\alpha=\kappa(p)$.

As for the exponent $d(p)$ of equation (50), the conventions $\sqrt t=0$ for $t<0$ and $1/0=0$ make
the formula total; the paper uses it only for $1<p<\infty$, where $p-1>0$ and
the radicand is positive.
\end{definition}
